%% Abstract (configurado para língua inglesa)
\begin{resumo}[Abstract]
\begin{otherlanguage*}{english}

The increase in the number of people living alone, especially older adults and individuals with special needs, has raised concerns regarding emergency situations occurring without immediate assistance. In this context, this work presents the development of Linha Vital, a mobile system designed to monitor periods of inactivity and automatically notify trusted contacts in risk situations. The objective of the project is to provide an accessible, non-invasive technological solution capable of contributing to the safety and autonomy of vulnerable individuals. The adopted methodology involved a literature review on remote health monitoring, requirements elicitation, system modeling, definition of business rules, development of user stories, and design of the application architecture. The solution was developed using Kotlin for both mobile and backend development, Spring Boot for service implementation, PostgreSQL for data persistence, and Oracle Cloud Infrastructure for hosting the application components. The system operates by identifying prolonged periods of inactivity, sending periodic well-being verification notifications, and automatically triggering alerts to previously registered contacts when risk conditions are detected. In addition, it provides manual emergency activation features and accessibility-oriented mechanisms. The results demonstrate the technical feasibility of the proposed solution, showing that the use of consolidated technologies enables the development of a reliable, scalable, and user-centered system. It is concluded that Linha Vital represents a promising alternative for improving the safety of individuals living alone, contributing to reduced response times in emergency situations and promoting a better quality of life through the integration of technology, accessibility, and care.

\vspace{\onelineskip}
\noindent
\textbf{Keywords}: remote monitoring; assistive technology; accessibility; mobile application; digital health.

\end{otherlanguage*}
\end{resumo}


%% Exemplo de resumo em francês
%\begin{resumo}[Résumé]
% \begin{otherlanguage*}{french}
%    Il s'agit d'un résumé en français.
% 
%   \textbf{Mots-clés}: latex. abntex. publication de textes.
% \end{otherlanguage*}
%\end{resumo}

%% Exemplo de resumo em Espanhol
%\begin{resumo}[Resumen]
% \begin{otherlanguage*}{spanish}
%   Este es el resumen en español.
%  
%   \textbf{Palabras clave}: latex. abntex. publicación de textos.
% \end{otherlanguage*}
%\end{resumo}
% ---